% Options for packages loaded elsewhere
\PassOptionsToPackage{unicode}{hyperref}
\PassOptionsToPackage{hyphens}{url}
\PassOptionsToPackage{dvipsnames,svgnames,x11names}{xcolor}
\PassOptionsToPackage{space}{xeCJK}
%
\documentclass[
  a4paper,
  12pt]{ctexart}
\usepackage{amsmath,amssymb}
\usepackage{iftex}
\ifPDFTeX
  \usepackage[T1]{fontenc}
  \usepackage[utf8]{inputenc}
  \usepackage{textcomp} % provide euro and other symbols
\else % if luatex or xetex
  \usepackage{unicode-math} % this also loads fontspec
  \defaultfontfeatures{Scale=MatchLowercase}
  \defaultfontfeatures[\rmfamily]{Ligatures=TeX,Scale=1}
\fi
\usepackage{lmodern}
\ifPDFTeX\else
  % xetex/luatex font selection
    \setmainfont[]{Noto Serif}
    \setsansfont[]{Noto Sans}
    \setmonofont[]{DejaVu Sans Mono}
  \ifXeTeX
    \usepackage{xeCJK}
    \setCJKmainfont[]{Noto Serif CJK SC}
          \setCJKsansfont[]{Noto Sans CJK SC}
          \fi
  \ifLuaTeX
    \usepackage[]{luatexja-fontspec}
    \setmainjfont[]{Noto Serif CJK SC}
  \fi
\fi
% Use upquote if available, for straight quotes in verbatim environments
\IfFileExists{upquote.sty}{\usepackage{upquote}}{}
\IfFileExists{microtype.sty}{% use microtype if available
  \usepackage[]{microtype}
  \UseMicrotypeSet[protrusion]{basicmath} % disable protrusion for tt fonts
}{}
\makeatletter
\@ifundefined{KOMAClassName}{% if non-KOMA class
  \IfFileExists{parskip.sty}{%
    \usepackage{parskip}
  }{% else
    \setlength{\parindent}{0pt}
    \setlength{\parskip}{6pt plus 2pt minus 1pt}}
}{% if KOMA class
  \KOMAoptions{parskip=half}}
\makeatother
\usepackage{xcolor}
\usepackage[margin=22mm]{geometry}
\usepackage{color}
\usepackage{fancyvrb}
\newcommand{\VerbBar}{|}
\newcommand{\VERB}{\Verb[commandchars=\\\{\}]}
\DefineVerbatimEnvironment{Highlighting}{Verbatim}{commandchars=\\\{\}}
% Add ',fontsize=\small' for more characters per line
\newenvironment{Shaded}{}{}
\newcommand{\AlertTok}[1]{\textcolor[rgb]{1.00,0.00,0.00}{\textbf{#1}}}
\newcommand{\AnnotationTok}[1]{\textcolor[rgb]{0.38,0.63,0.69}{\textbf{\textit{#1}}}}
\newcommand{\AttributeTok}[1]{\textcolor[rgb]{0.49,0.56,0.16}{#1}}
\newcommand{\BaseNTok}[1]{\textcolor[rgb]{0.25,0.63,0.44}{#1}}
\newcommand{\BuiltInTok}[1]{\textcolor[rgb]{0.00,0.50,0.00}{#1}}
\newcommand{\CharTok}[1]{\textcolor[rgb]{0.25,0.44,0.63}{#1}}
\newcommand{\CommentTok}[1]{\textcolor[rgb]{0.38,0.63,0.69}{\textit{#1}}}
\newcommand{\CommentVarTok}[1]{\textcolor[rgb]{0.38,0.63,0.69}{\textbf{\textit{#1}}}}
\newcommand{\ConstantTok}[1]{\textcolor[rgb]{0.53,0.00,0.00}{#1}}
\newcommand{\ControlFlowTok}[1]{\textcolor[rgb]{0.00,0.44,0.13}{\textbf{#1}}}
\newcommand{\DataTypeTok}[1]{\textcolor[rgb]{0.56,0.13,0.00}{#1}}
\newcommand{\DecValTok}[1]{\textcolor[rgb]{0.25,0.63,0.44}{#1}}
\newcommand{\DocumentationTok}[1]{\textcolor[rgb]{0.73,0.13,0.13}{\textit{#1}}}
\newcommand{\ErrorTok}[1]{\textcolor[rgb]{1.00,0.00,0.00}{\textbf{#1}}}
\newcommand{\ExtensionTok}[1]{#1}
\newcommand{\FloatTok}[1]{\textcolor[rgb]{0.25,0.63,0.44}{#1}}
\newcommand{\FunctionTok}[1]{\textcolor[rgb]{0.02,0.16,0.49}{#1}}
\newcommand{\ImportTok}[1]{\textcolor[rgb]{0.00,0.50,0.00}{\textbf{#1}}}
\newcommand{\InformationTok}[1]{\textcolor[rgb]{0.38,0.63,0.69}{\textbf{\textit{#1}}}}
\newcommand{\KeywordTok}[1]{\textcolor[rgb]{0.00,0.44,0.13}{\textbf{#1}}}
\newcommand{\NormalTok}[1]{#1}
\newcommand{\OperatorTok}[1]{\textcolor[rgb]{0.40,0.40,0.40}{#1}}
\newcommand{\OtherTok}[1]{\textcolor[rgb]{0.00,0.44,0.13}{#1}}
\newcommand{\PreprocessorTok}[1]{\textcolor[rgb]{0.74,0.48,0.00}{#1}}
\newcommand{\RegionMarkerTok}[1]{#1}
\newcommand{\SpecialCharTok}[1]{\textcolor[rgb]{0.25,0.44,0.63}{#1}}
\newcommand{\SpecialStringTok}[1]{\textcolor[rgb]{0.73,0.40,0.53}{#1}}
\newcommand{\StringTok}[1]{\textcolor[rgb]{0.25,0.44,0.63}{#1}}
\newcommand{\VariableTok}[1]{\textcolor[rgb]{0.10,0.09,0.49}{#1}}
\newcommand{\VerbatimStringTok}[1]{\textcolor[rgb]{0.25,0.44,0.63}{#1}}
\newcommand{\WarningTok}[1]{\textcolor[rgb]{0.38,0.63,0.69}{\textbf{\textit{#1}}}}
\setlength{\emergencystretch}{3em} % prevent overfull lines
\providecommand{\tightlist}{%
  \setlength{\itemsep}{0pt}\setlength{\parskip}{0pt}}
\setcounter{secnumdepth}{-\maxdimen} % remove section numbering
\ifLuaTeX
  \usepackage{selnolig}  % disable illegal ligatures
\fi
\usepackage{bookmark}
\IfFileExists{xurl.sty}{\usepackage{xurl}}{} % add URL line breaks if available
\urlstyle{same}
\hypersetup{
  colorlinks=true,
  linkcolor={black},
  filecolor={Maroon},
  citecolor={Blue},
  urlcolor={blue},
  pdfcreator={LaTeX via pandoc}}

\author{}
\date{}

\begin{document}

{
\hypersetup{linkcolor=}
\setcounter{tocdepth}{2}
\tableofcontents
}
\section{Sondow-Pudlak
程序中的一个参数闭合形式化碰撞定理}\label{sondow-pudlak-ux7a0bux5e8fux4e2dux7684ux4e00ux4e2aux53c2ux6570ux95edux5408ux5f62ux5f0fux5316ux78b0ux649eux5b9aux7406}

作者：James

\subsection{摘要}\label{ux6458ux8981}

本文给出一个 Lean 4 形式化的 Sondow-Pudlak 型碰撞定理。设
\texttt{minCheckedCodeSize\ n} 表示同一 PA-Hilbert
检查器坐标中的最小已检查证明码 长度。本文的主定理把旧路线中显式出现的
\texttt{upper\_provider}、\texttt{tail\_gap} 和
\texttt{eventually\_strict\_length}
全部从定理表面移除，剩余输入压缩为一个清楚的
单项式增长下界：对任意自然数 \texttt{coeff,\ degree}，存在显式阈值
\texttt{thresholdOfMonomial\ coeff\ degree}，使得阈值以后

\begin{Shaded}
\begin{Highlighting}[]
\NormalTok{coeff * (n + 1)\^{}degree \textless{} minCheckedCodeSize n.}
\end{Highlighting}
\end{Shaded}

在该输入下，Lean 证明碰撞指标具有公式

\begin{Shaded}
\begin{Highlighting}[]
\NormalTok{bigN = thresholdOfMonomial upper\_data.coeff upper\_data.degree}
\end{Highlighting}
\end{Shaded}

并在同一指标处推出上界侧与检查器侧的形式化矛盾，从而得到
\texttt{¬\ is\_rational\ euler\_mascheroni}。主定理名称、复现命令和
residual-constant 审计见第 4、5 节；观察到的 axiom profile 只含标准
Lean/Mathlib 逻辑依赖。

关键词：Euler-Mascheroni 常数；Sondow 判据；Pudlak 下界；证明长度；
形式化数学；Lean。

\subsection{1. 引言}\label{ux5f15ux8a00}

Euler-Mascheroni 常数

\begin{Shaded}
\begin{Highlighting}[]
\NormalTok{\textbackslash{}gamma=\textbackslash{}lim\_\{n\textbackslash{}to\textbackslash{}infty\}\textbackslash{}left(\textbackslash{}sum\_\{k=1\}\^{}\{n\}\textbackslash{}frac\{1\}\{k\}{-}\textbackslash{}log n\textbackslash{}right)}
\end{Highlighting}
\end{Shaded}

的无理性是经典困难问题。Sondow
的判据从有理性假设出发产生可计算的算术上界 结构；Pudlak
型证明复杂度下界则从证明长度方向给出反向限制。若两侧能够在
同一个检查器测度坐标中相遇，便可在某个自然数处形成

\begin{Shaded}
\begin{Highlighting}[]
\NormalTok{M(N)\textbackslash{}le U(N)\textless{}M(N)}
\end{Highlighting}
\end{Shaded}

式的碰撞。

本文的贡献是把这一碰撞机制整理成一个参数表面更干净的 Lean 定理。旧的
中间路线把 \texttt{upper\_provider} 与 \texttt{tail\_gap}
直接暴露在主定理中；本文使用
项目长度端点，把它们上移并合并到更具体的单项式增长输入中。这样，审稿人需要
检查的核心数学义务不再是不透明的证书对象，而是如下明确命题：

\begin{Shaded}
\begin{Highlighting}[]
\NormalTok{T(c,d)\textbackslash{}le n\textbackslash{}quad\textbackslash{}Longrightarrow\textbackslash{}quad}
\NormalTok{c(n+1)\^{}d\textless{}\textbackslash{}operatorname\{minCheckedCodeSize\}(n).}
\end{Highlighting}
\end{Shaded}

这里 \texttt{T(c,d)} 即 Lean 中的 \texttt{thresholdOfMonomial\ c\ d}。

本文不声称已经给出十进制数值 \texttt{N}，也不把旧 half-denominator
精细公式作为 本稿主结果。本文主张的是：在上述明确增长输入下，Lean
已经验证了公式级 碰撞指标、同点矛盾以及干净 axiom profile。

\subsection{2. 形式化对象}\label{ux5f62ux5f0fux5316ux5bf9ux8c61}

固定一个 PA-Hilbert 检查器坐标。对左右两个形式化证明族
\texttt{left\_family} 和 \texttt{right\_family}，Lean
构造合取引入后再右合取消去的证明族，并定义

\begin{Shaded}
\begin{Highlighting}[]
\NormalTok{M n =}
\NormalTok{  ((left\_family.conjIntro right\_family)}
\NormalTok{    |\textgreater{}.rightConjElim}
\NormalTok{    |\textgreater{}.minCheckedCodeSize n)}
\end{Highlighting}
\end{Shaded}

这就是本文的 \texttt{minCheckedCodeSize\ n}。

Sondow 侧的上界数据被压缩为一个自然数幂上界数据

\begin{Shaded}
\begin{Highlighting}[]
\NormalTok{upper\_data =}
\NormalTok{  projectLengthConjIntroLengthAddTwoNatPowerUpperData}
\NormalTok{    left\_family right\_family left\_data right\_data}
\end{Highlighting}
\end{Shaded}

其中 \texttt{upper\_data.coeff} 和 \texttt{upper\_data.degree}
给出生成上界的自然数单项式 包络。于是本文的碰撞指标定义为

\begin{Shaded}
\begin{Highlighting}[]
\NormalTok{bigN =}
\NormalTok{  thresholdOfMonomial upper\_data.coeff upper\_data.degree.}
\end{Highlighting}
\end{Shaded}

主输入是单项式增长证书：

\begin{Shaded}
\begin{Highlighting}[]
\NormalTok{monomial\_lt\_lengthCodeAt\_after :}
\NormalTok{  ∀ coeff degree n : Nat,}
\NormalTok{    thresholdOfMonomial coeff degree ≤ n →}
\NormalTok{      coeff * (n + 1)\^{}degree \textless{} M n}
\end{Highlighting}
\end{Shaded}

它是当前定理的真正数学承重部分。

\subsection{3. 推导}\label{ux63a8ux5bfc}

令

\begin{Shaded}
\begin{Highlighting}[]
\NormalTok{c=\textbackslash{}texttt\{upper\textbackslash{}\_data.coeff\},\textbackslash{}qquad}
\NormalTok{d=\textbackslash{}texttt\{upper\textbackslash{}\_data.degree\},\textbackslash{}qquad}
\NormalTok{N=T(c,d).}
\end{Highlighting}
\end{Shaded}

由单项式增长输入，因 \texttt{T(c,d)\ ≤\ N}，得到

\begin{Shaded}
\begin{Highlighting}[]
\NormalTok{c(N+1)\^{}d\textless{}M(N).}
\end{Highlighting}
\end{Shaded}

另一方面，Sondow 侧的形式化上界构造保证生成的上界函数在 \texttt{N}
处被同一个 单项式包络控制。用 Lean 端点的记号写，就是

\begin{Shaded}
\begin{Highlighting}[]
\NormalTok{U(N)\textbackslash{}le c(N+1)\^{}d.}
\end{Highlighting}
\end{Shaded}

合并二者得到

\begin{Shaded}
\begin{Highlighting}[]
\NormalTok{U(N)\textless{}M(N).}
\end{Highlighting}
\end{Shaded}

同时，检查器端点把同一个 \texttt{N}
放回到共同测度坐标中，得到反向的已检查上界

\begin{Shaded}
\begin{Highlighting}[]
\NormalTok{M(N)\textbackslash{}le U(N).}
\end{Highlighting}
\end{Shaded}

于是

\begin{Shaded}
\begin{Highlighting}[]
\NormalTok{M(N)\textbackslash{}le U(N)\textless{}M(N),}
\end{Highlighting}
\end{Shaded}

矛盾。Lean 中的主证书还记录了 \texttt{N} 的程序化公式：

\begin{Shaded}
\begin{Highlighting}[]
\NormalTok{bigN = thresholdOfMonomial upper\_data.coeff upper\_data.degree.}
\end{Highlighting}
\end{Shaded}

这正是本文保留的公式级大 \texttt{N}。

\subsection{4. 主定理}\label{ux4e3bux5b9aux7406}

主定理 Lean 名称为

\begin{Shaded}
\begin{Highlighting}[]
\NormalTok{singletonMonomialLowerBound\_submissionRoute}
\end{Highlighting}
\end{Shaded}

源码：

\url{https://github.com/james130012/sondow-pudlak-collision-box/blob/main/integration/SondowProjectBigNParameterClosureAudit.lean}

省略 namespace 后，其数学内容可写为：

\begin{Shaded}
\begin{Highlighting}[]
\NormalTok{theorem singletonMonomialLowerBound\_submissionRoute}
\NormalTok{    {-}{-} other explicit arguments are shown in the source file}
\NormalTok{    (thresholdOfMonomial : Nat → Nat → Nat)}
\NormalTok{    (monomial\_lt\_lengthCodeAt\_after :}
\NormalTok{      ∀ coeff degree n : Nat,}
\NormalTok{        thresholdOfMonomial coeff degree ≤ n →}
\NormalTok{          coeff * (n + 1)\^{}degree \textless{} minCheckedCodeSize n) :}
\NormalTok{    bigN = thresholdOfMonomial upper\_data.coeff upper\_data.degree ∧}
\NormalTok{      ¬ is\_rational euler\_mascheroni}
\end{Highlighting}
\end{Shaded}

注释行所指的显式参数包括左右证明族、严格时间界、自然数幂上界数据以及这些长度函数
为多项式有界的 Lean 证书。它们都是显式 Lean 参数；主定理表面不再出现
\texttt{upper\_provider}、\texttt{tail\_gap} 或
\texttt{eventually\_strict\_length}。

\subsection{5. 形式化审计}\label{ux5f62ux5f0fux5316ux5ba1ux8ba1}

复现命令如下：

\begin{Shaded}
\begin{Highlighting}[]
\ExtensionTok{lake}\NormalTok{ exe cache get}
\ExtensionTok{lake}\NormalTok{ build integration.SondowProjectBigNParameterClosureAudit}
\ExtensionTok{lake}\NormalTok{ env lean }\AttributeTok{{-}{-}stdin} \OperatorTok{\textless{}\textless{}\textquotesingle{}EOF\textquotesingle{}}
\StringTok{import integration.SondowProjectBigNParameterClosureAudit}
\StringTok{open SondowMainCheckedCodeBridge.SondowProjectBigNParameterClosureAudit}

\StringTok{\#check singletonMonomialLowerBound\_submissionRoute}
\StringTok{\#print axioms singletonMonomialLowerBound\_submissionRoute}
\OperatorTok{EOF}
\end{Highlighting}
\end{Shaded}

观察到的输出为：

\begin{Shaded}
\begin{Highlighting}[]
\NormalTok{[propext, Classical.choice, Quot.sound]}
\end{Highlighting}
\end{Shaded}

未出现以下项目级 residual constants：

\begin{Shaded}
\begin{Highlighting}[]
\NormalTok{partial\_consistency\_payload}
\NormalTok{proof\_length}
\NormalTok{strengthened\_partial\_consistency\_payload}
\end{Highlighting}
\end{Shaded}

审计报告：

\begin{itemize}
\tightlist
\item
  中文：\url{https://github.com/james130012/sondow-pudlak-collision-box/blob/main/docs/parameter_closure_audit_20260708_zh.md}
\item
  English：\url{https://github.com/james130012/sondow-pudlak-collision-box/blob/main/docs/parameter_closure_audit_20260708_en.md}
\end{itemize}

仓库：

\url{https://github.com/james130012/sondow-pudlak-collision-box}

\subsection{6. 结论}\label{ux7ed3ux8bba}

本文给出一个参数表面闭合的 Lean 形式化碰撞定理。旧路线中的
\texttt{upper\_provider}、\texttt{tail\_gap} 和
\texttt{eventually\_strict\_length} 不再作为主定理
表面参数出现；当前主定理只暴露一个清楚的单项式增长下界，并在该输入下证明
公式级大 \texttt{N} 与同点矛盾。

后续工作是关闭该单项式增长下界本身，并在需要时继续构造 half-denominator
精细公式和十进制数值阈值。

\subsection{参考文献}\label{ux53c2ux8003ux6587ux732e}

\begin{enumerate}
\def\labelenumi{\arabic{enumi}.}
\tightlist
\item
  J. Sondow, Criteria for irrationality of Euler's constant,
  \emph{Proceedings of the American Mathematical Society} 131 (2003),
  3335-3344.
\item
  S. R. Buss, On Godel's theorems on lengths of proofs I: Number of
  lines and speedup for arithmetics, \emph{Journal of Symbolic Logic} 59
  (1994), 737-756.
\item
  P. Pudlak, On the lengths of proofs of finitistic consistency
  statements in first order theories, in \emph{Logic Colloquium 1984},
  North-Holland, 1986.
\item
  L. de Moura and S. Ullrich, The Lean 4 theorem prover and programming
  language, in \emph{Automated Deduction - CADE 28}, 2021.
\end{enumerate}

\end{document}
